% Chapter — Engine propulsion with thrust-vector control (FR-10). Follows
% the mandatory FR-29 six-section template: purpose; assumptions and
% domain of validity; derivation; discretization and implementation;
% validation evidence; references. The equation labels
% eq:propulsion:thrust, eq:propulsion:mdot, eq:propulsion:direction,
% eq:propulsion:forcetorque, eq:propulsion:tsiolkovsky,
% eq:propulsion:spool, and eq:propulsion:gimbalslew are echoed verbatim in
% comments in cpp/src/models/propulsion.cpp (FR-29
% equation-label-to-code traceability); renaming them breaks that trace.
\chapter{Engine Propulsion and Thrust-Vector Control}
\label{ch:propulsion}

\section{Purpose}
\label{sec:propulsion:purpose}

This chapter documents the FR-10 engine model: per-engine delivered
thrust with ambient back-pressure, mass flow from the constant vacuum
specific impulse, throttle limits, linear spool-up and spool-down ramps,
ignition-count enforcement, and gimbal (TVC) angle and rate limits. The
engine produces a body force and a torque about the current composite CG
(Chapter~\ref{ch:massprops}) and a propellant mass-flow rate that drives
the settled-tank depletion. The model is implemented in
\texttt{cpp/src/models/propulsion.cpp} with its interface in
\texttt{cpp/include/star/models/propulsion.hpp}.

\section{Assumptions and domain of validity}
\label{sec:propulsion:domain}

Assumptions:
\begin{enumerate}
  \item \textbf{Constant vacuum specific impulse.} Each engine's
    $I_{\mathrm{sp,vac}}$ is constant over its throttle range (FR-10):
    chamber conditions scale proportionally with throttle, so vacuum
    thrust and mass flow scale together. Real engines lose a few
    seconds of $I_{\mathrm{sp}}$ at deep throttle; that error is
    absorbed into the configured rating.
  \item \textbf{Back pressure acts only on delivered thrust.} The
    ambient term $p_{\mathrm{amb}} A_e$ is the static-pressure force on
    the nozzle exit plane; it does not alter the internal flow, so the
    mass flow follows the vacuum rating and throttle only
    (equation~\eqref{eq:propulsion:mdot}). Flow separation in heavily
    overexpanded operation is not modeled; the delivered thrust may go
    negative if an engine is run at low throttle deep in an atmosphere,
    and the configuration layer is responsible for physically sensible
    throttle floors.
  \item \textbf{Rigid TVC geometry.} The engine is mounted at a fixed
    position with a fixed nominal thrust axis in the vehicle frame
    (FR-13, $+X$ forward); gimbal deflection rotates the thrust
    direction but not the application point. Inertial (tail-wags-dog)
    effects of the swiveling engine mass are not modeled.
  \item \textbf{Quasi-static throttle response.} Spool-up and spool-down
    are linear ramps of the delivered throttle level between 0 and the
    commanded setting (FR-10); combustion transients are not modeled.
\end{enumerate}

Domain of validity: non-negative vacuum thrust and exit area; positive
$I_{\mathrm{sp,vac}}$; throttle limits
$0 \le \lambda_{\min} \le \lambda_{\max} \le 1$; unit-norm nominal axis
and mutually orthogonal unit gimbal axes; non-negative spool times,
gimbal limit, and gimbal rate; ambient pressure $p_{\mathrm{amb}} \ge 0$
supplied by the atmosphere model in use (USSA76 for Earth ascent,
Chapter~\ref{ch:ussa76}). Out-of-domain behavior, matching the code
exactly: non-finite inputs, negative ratings, malformed axis vectors,
inverted throttle limits, and negative step sizes throw
\texttt{std::domain\_error}. A zero delivered throttle level returns
exactly zero force, torque, and mass flow --- an engine that is off
exerts no back-pressure force because its nozzle is not flowing.

\section{Derivation}
\label{sec:propulsion:derivation}

\subsection{Delivered thrust and mass flow}

Consider a steady control volume enclosing the engine, cut at the
nozzle exit plane of area $A_e$. Momentum conservation for exhaust
leaving at effective exhaust speed $v_e$ with mass flow $\dot{m}$ and
exit static pressure $p_e$, immersed in ambient pressure
$p_{\mathrm{amb}}$ (which acts on every outer surface of the vehicle
and leaves only the exit-plane imbalance uncancelled), gives the
delivered thrust
\begin{equation*}
  F \;=\; \dot{m} v_e + (p_e - p_{\mathrm{amb}})\, A_e.
\end{equation*}
Define the \emph{vacuum thrust} at the same operating point,
$F_{\mathrm{vac}} = \dot{m} v_e + p_e A_e$, i.e.\ the thrust the same
chamber state delivers with $p_{\mathrm{amb}} = 0$. With a throttle
level $\lambda$ scaling the chamber mass flow (and hence
$F_{\mathrm{vac}}$) proportionally,
\begin{equation}
  \boxed{\; F \;=\; \lambda\, F_{\mathrm{vac}}
    - p_{\mathrm{amb}}\, A_e \quad (\lambda > 0); \qquad
    F \;=\; 0 \quad (\lambda = 0). \;}
  \label{eq:propulsion:thrust}
\end{equation}
The vacuum specific impulse is defined by
$F_{\mathrm{vac}} = \dot{m}\, g_0\, I_{\mathrm{sp,vac}}$ with the
standard gravity $g_0 = \qty{9.80665}{\metre\per\second\squared}$ exact
(the conventional value fixed by the 3rd CGPM, 1901, as published in the
SI Brochure~\cite{bipm2019}). Solving for the mass flow at throttle
$\lambda$:
\begin{equation}
  \boxed{\; \dot{m} \;=\;
    \frac{\lambda\, F_{\mathrm{vac}}}{g_0\, I_{\mathrm{sp,vac}}}. \;}
  \label{eq:propulsion:mdot}
\end{equation}
The mass flow is set by the vacuum rating and the throttle alone:
ambient back pressure reduces the \emph{delivered} thrust
\eqref{eq:propulsion:thrust} but never the propellant consumption
\eqref{eq:propulsion:mdot}. Equivalently, the delivered specific
impulse $F / (\dot{m} g_0)$ falls with altitude while
$I_{\mathrm{sp,vac}}$ stays constant.

\subsection{Gimbal deflection and force/torque coupling}

The engine's nominal thrust direction is the unit vector $\hat{\vv{u}}$
(vehicle frame; $+X$ for an aft-mounted main engine). TVC deflects it
by two gimbal angles $(\gamma_1, \gamma_2)$ about two configured
orthogonal unit axes $\hat{\vv{a}}_1, \hat{\vv{a}}_2$ perpendicular to
$\hat{\vv{u}}$ (defaults $+Y$, $+Z$), applied in that fixed order with
the Rodrigues rotation
$\mat{R}(\hat{\vv{a}}, \gamma)\vv{v} = \vv{v}\cos\gamma +
(\hat{\vv{a}} \times \vv{v})\sin\gamma + \hat{\vv{a}}
(\hat{\vv{a}} \cdot \vv{v})(1 - \cos\gamma)$:
\begin{equation}
  \hat{\vv{d}} \;=\; \mat{R}(\hat{\vv{a}}_2, \gamma_2)\,
    \mat{R}(\hat{\vv{a}}_1, \gamma_1)\, \hat{\vv{u}}.
  \label{eq:propulsion:direction}
\end{equation}
Both factors are proper rotations, so $\hat{\vv{d}}$ is exactly unit
norm up to rounding. With the engine mounted at $\vv{r}_e$ and the
composite CG at $\bar{\vv{c}}$, the body force and the torque about the
CG are
\begin{equation}
  \boxed{\; \vv{F} \;=\; F\, \hat{\vv{d}}, \qquad
  \vv{\tau} \;=\; (\vv{r}_e - \bar{\vv{c}}) \times \vv{F}. \;}
  \label{eq:propulsion:forcetorque}
\end{equation}

\subsection{Tsiolkovsky benchmark}

For a fixed-attitude vacuum burn at constant throttle, the
translational speed obeys $\mathrm{d}v/\mathrm{d}t = \lambda
F_{\mathrm{vac}} / m(t)$ with $m(t) = m_0 - \dot{m} t$ and $\dot{m}$
from \eqref{eq:propulsion:mdot}. Substituting
$\lambda F_{\mathrm{vac}} = \dot{m}\, g_0 I_{\mathrm{sp,vac}}$ and
integrating from $m_0$ to $m_1 = m(t_{\mathrm{burn}})$:
\begin{equation}
  \Delta v \;=\; \int_0^{t_b}
    \frac{\dot{m}\, g_0 I_{\mathrm{sp,vac}}}{m_0 - \dot{m} t}\,
    \mathrm{d}t
  \;=\; g_0\, I_{\mathrm{sp,vac}} \ln \frac{m_0}{m_1},
  \label{eq:propulsion:tsiolkovsky}
\end{equation}
the rocket equation. It is exact under this chapter's assumptions and
independent of the throttle profile (only the endpoint masses enter),
which makes it the Phase 4 exit-criterion-3 acceptance instrument for
the thrust/mass-flow bookkeeping.

\section{Discretization and implementation}
\label{sec:propulsion:implementation}

\begin{enumerate}
  \item \textbf{Module and traceability.} The model lives in
    \texttt{cpp/src/models/propulsion.cpp}; the source comments echo
    the labels \texttt{eq:propulsion:thrust} through
    \texttt{eq:propulsion:gimbalslew} verbatim.
  \item \textbf{State advance vs.\ pure evaluation.} The actuator state
    (running flag, ignition count, delivered throttle level, gimbal
    angles) advances once per step by a discrete update; force, torque,
    and mass flow are pure functions of parameters and state so the
    integrator may evaluate them at any stage without side effects
    (D-2 core discipline: plain structs in, plain structs out).
  \item \textbf{Spool ramp.} The delivered level $\lambda$ slews
    linearly toward the target
    $\lambda^\ast = \operatorname{clamp}(\lambda_{\mathrm{cmd}},
    \lambda_{\min}, \lambda_{\max})$ while running and toward $0$ when
    shut down, at the configured full-range rates:
    \begin{equation}
      \lambda' \;=\; \lambda + \operatorname{clamp}\!\bigl(
        \lambda^\ast - \lambda,\;
        -\Delta t / t_{\mathrm{down}},\;
        +\Delta t / t_{\mathrm{up}} \bigr),
      \label{eq:propulsion:spool}
    \end{equation}
    where $t_{\mathrm{up}}$ ($t_{\mathrm{down}}$) is the 0-to-1
    (1-to-0) spool time; a zero spool time applies the target
    immediately. The ramp is exactly linear: away from the endpoints
    each step adds exactly $\Delta t / t_{\mathrm{up}}$.
  \item \textbf{Ignition count.} A run command on a non-running engine
    consumes one ignition; when the configured count is exhausted the
    command is refused and the engine stays off (the level target
    remains $0$). Shutting down mid-spool and re-igniting consumes a
    new count.
  \item \textbf{Gimbal slew and clamp.} Each gimbal angle slews toward
    its command at exactly the configured rate and is clamped at the
    angle limit, per axis:
    \begin{equation}
      \gamma_i' \;=\; \operatorname{clamp}\!\bigl(
        \gamma_i + \operatorname{clamp}(
          \gamma_{\mathrm{cmd},i} - \gamma_i,\;
          -\dot{\gamma}_{\max} \Delta t,\;
          +\dot{\gamma}_{\max} \Delta t ),\;
        -\gamma_{\max},\; +\gamma_{\max} \bigr).
      \label{eq:propulsion:gimbalslew}
    \end{equation}
    A step command therefore produces a ramp whose per-step increment
    is exactly $\dot{\gamma}_{\max} \Delta t$ until the final partial
    step (Phase 4 exit criterion 7). The angle limit is a per-axis
    (square) clamp, not a cone.
  \item \textbf{Determinism.} The only libm calls are the
    $\sin/\cos$ pair per gimbal axis in
    \eqref{eq:propulsion:direction}; everything else is IEEE-754 basic
    operations in fixed order. Zero throttle returns literal zeros.
\end{enumerate}

\section{Validation evidence}
\label{sec:propulsion:validation}

Table~\ref{tab:propulsion:validation} lists the tests that constitute
this chapter's validation evidence. Each test identifier is
machine-checked to exist in the test tree by
\texttt{scripts/lint\_chapters.py}; golden provenance and tolerances
are recorded in \texttt{tests/golden/propulsion/manifest.toml}.

\begin{table}[htbp]
  \centering
  \caption{Validation evidence for the propulsion model (doctest,
    \texttt{cpp/tests/test\_propulsion.cpp}).}
  \label{tab:propulsion:validation}
  \begin{tabular}{lp{8.6cm}}
    \toprule
    Test ID & Acceptance basis \\
    \midrule
    \testid{PROP-ENGINE-GOLDEN} &
      Five committed operating points (vacuum, sea level, mid-throttle,
      gimballed, zero-throttle) match the independent 60-digit mpmath
      evaluation of \eqref{eq:propulsion:thrust},
      \eqref{eq:propulsion:mdot}, \eqref{eq:propulsion:direction}, and
      \eqref{eq:propulsion:forcetorque} to $10^{-12}$ relative
      (norm-relative for vectors); the zero-throttle case is exactly
      zero in every output. \\
    \testid{PROP-BACKPRESSURE-USSA76} &
      The sea-level operating point evaluated with
      $p_{\mathrm{amb}}$ taken live from the USSA76 model at $z = 0$
      (Chapter~\ref{ch:ussa76}) reproduces the committed golden thrust:
      the exit-criterion-3 back-pressure anchor
      $F = F_{\mathrm{vac}} - \qty{101325}{\pascal} \cdot A_e$. \\
    \testid{PROP-TSIOLKOVSKY-BURN} &
      RK4 integration of the fixed-attitude vacuum burn through the
      model's thrust and mass flow accumulates a $\Delta v$ matching
      \eqref{eq:propulsion:tsiolkovsky} within $0.1\,\%$ (Phase 4 exit
      criterion 3), with the endpoint mass checked at $10^{-12}$
      relative. \\
    \testid{PROP-SPOOL-RAMP} &
      The delivered level under \eqref{eq:propulsion:spool} rises by
      exactly $\Delta t / t_{\mathrm{up}}$ per step from ignition to
      the throttle target, holds exactly, and falls by exactly
      $\Delta t / t_{\mathrm{down}}$ per step after cutoff to exactly
      zero. \\
    \testid{PROP-THROTTLE-CLAMP} &
      Commands outside $[\lambda_{\min}, \lambda_{\max}]$ settle
      exactly on the violated bound; in-range commands settle exactly
      on the command. \\
    \testid{PROP-IGNITION-COUNT} &
      An engine with $N$ configured ignitions accepts exactly $N$
      run transitions and refuses the $(N{+}1)$-th, remaining at zero
      level with zero force and mass flow. \\
    \testid{PROP-TVC-RATE-LIMIT} &
      A step gimbal command produces per-step increments of exactly
      $\dot{\gamma}_{\max} \Delta t$ (Phase 4 exit criterion 7) and the
      angle clamps exactly at $\pm\gamma_{\max}$ under an
      over-limit command, per \eqref{eq:propulsion:gimbalslew}. \\
    \testid{PROP-DOMAIN} &
      Out-of-domain inputs per Section~\ref{sec:propulsion:domain}
      throw \texttt{std::domain\_error}. \\
    \bottomrule
  \end{tabular}
\end{table}

\section{References}
\label{sec:propulsion:references}

The thrust and mass-flow relations, the gimbal geometry, and the
Tsiolkovsky benchmark are derived from first principles in
Section~\ref{sec:propulsion:derivation} (momentum balance over a
steady nozzle control volume and the definition of vacuum specific
impulse); no listed project reference covers the rocket-propulsion
formulation, so the derivation is self-contained by design. The exact
conventional standard gravity $g_0 = \qty{9.80665}{\metre\per\second
\squared}$ is taken from the BIPM SI Brochure~\cite{bipm2019}; the
sea-level anchor pressure is the U.S. Standard Atmosphere, 1976
standard sea-level pressure~\cite{ussa1976} as served by
Chapter~\ref{ch:ussa76}.
